\chapter{Infinite-dimensional symplectic reduction} \label{Chapter_inf_dim_sympl_red}
The aim of this chapter is to give a Banach-manifold version of symplectic reduction (theorem \ref{Th_Marsden_Meyer_Weinstein}), for the action of a Banach Lie group on a Banach manifold. We refer the reader to appendix \ref{appendix_Banach_manifolds} for the relevant notions concerning infinite-dimensional Banach manifolds. The chapter starts with the (straighforward) generalization of symplectic manifolds to infinite-dimensiones. Once this is done, we can give the definition of moment map, and state the reduction theorem for Banach manifolds. As we shall see, the proof does not really change much from the classical theorem, aside from the fact that we no longer have a quotient manifold theorem, and that has to be addressed by requiring additional properties to moment maps. However, if we assume that we can define quotients, and they are Banach manifolds, the proof is quite similar to the finite-dimensional case. In fact, the original article by Marsden and Weinstein \parencite{MARSDEN1974121} gives a proof of theorem \ref{Th_Marsden_Meyer_Weinstein} in the context of Banach manifolds. A very thorough exposition of the theorem can be found in \parencite{Gianetto}. This results will be used in chapter \ref{chapter_HK_correspondence} to justify the fact that the Hitchin-Kobayashi may be seen as an infinite-dimensional version of the Kempf-Ness theorem introduced in chapter  \ref{chapter_kempf_Ness}. To that end, we devote the second half of the chapter to justifying that the gauge group introduced in section \ref{section_gauge_group} can be seen (with some subtelties) as a Hilbert Lie group.  

\section{Symplectic reduction theorem for Banach-manifolds}
The definition of an infinite-dimensional symplectic manifold is quite similar to the finite-dimensional one except for the fact that we no longer ask \(\Omega^\flat\) to be an isomorphism between \(T_m\mathcal{M}\cong T_m\mathcal{M}^*\), but to be injective. Observe that in the finite-dimensional setting this already implies that \(\Omega^\flat\) is an isomorphism because \(T_m\mathcal{M}\) and its dual are vector spaces of the same fonote dimension.

\Definition{\textbf{Banach symplectic manifold}}{Let \(\mathcal{M}\) be a Banach manifold, we say that a pair \((\mathcal{M}, \Omega)\) is a \textit{Banach symplectic manifold}\index{Banach!symplectic manifold} if \(\Omega\) is a closed two-form satisfying:
\begin{enumerate}
    \item For each \(m \in \mathcal{M}\), \(\Omega_m \colon T_m\mathcal{M} \times T_m\mathcal{M} \to \R\) is continuous.
    \item For each \(m \in \mathcal{M}\), \(\Omega_m\) is non-degenerate, that is, if \(u \in T_m\mathcal{M}\) and \(\Omega_m(u,v)=0\) for all \(v \in T_xX\), then \(u=0\).
\end{enumerate}
We say that \(\Omega\) is a symplectic form.
}
The notions of symplectomorphism, hamiltonian action of a Banach Lie group on a Banach manifold and moment map are defined just as in the finite-dimensional case (see section \ref{sec_symplectic_reduction}). Observe that given a Banach Lie group \(G\) and \(\xi \in \mathrm{Lie}(G)\), the infinitesimal action, i.e. the vector field \(X_{\xi}\) can be defined just as in \eqref{def_hat_X} because Banach Lie groups have an exponential map. The first major difference with respect to the finite setting is the fact that we do not have a direct generalization of the quotient manifold theorem. What ends up working is to ask more conditions to the moment map. Before dwelling in the problem of the lack of a quotient manifold theorem, let us consider the quotient \(\mu^{-1}(0)/G\). The fact that \(\mu\) is \(G\)-equivariant means that \(\mu^{-1}(0)\) is invariant under the \(G\)-action (same justification as \eqref{eq_mu_invariant_G_action}). Hence, \(G\) acts on \(\mu^{-1}(0)\), and the \textit{reduced space} \(\mathcal{M}_{\mathrm{red}}:=\mu^{-1}(0)/G\) is well-defined as a topological space. In the classical theorem, one could show that \(0\) was always a regular value of \(\mu\), and this implied that \(\mu^{-1}(0)\) was an embedded submanifold of the original manifold. Since we may not have that condition, we have to impose it on \(\mu\). This motivates the following definition.

\Definition{}{Let \(G\) be a Banach-Lie group acting symplectically on a Banach-manifold \(\mathcal{M}\), and let \(\mu \colon \mathcal{M} \to \mathfrak{g}^*\) be a moment map for the action.
\begin{enumerate}
    \item We say that \(0 \in \mathrm{Lie}(G)^*\) is a \textit{weakly regular value} of \(\mu\) if \(\mu^{-1}(0)\) is an embedded submanifold of \(\mathcal{M}\) and \(T_x(\mu^{-1}(0))=\ker \diff_x \mu\).
    \item We say that \(\mu\) \textit{satisfies the slice condition} if the action of \(G\) on \(\mathcal{M}\) is free, and for every \(x \in \mu^{-1}(0)\) there is a submanifold \(S_x \subset \mu^{-1}(0)\), called a \textit{slice} such that 
    \[
    T_x(\mu^{-1}(0))=T_xS_x\oplus T_x (G\cdot x).
    \]
\end{enumerate}
}
Aside from the motivation we just gave for the definition of weakly regular value, the condition \(T_x(\mu^{-1}(0))=\ker \diff_x \mu\) is always satisfied for a moment map in the classical sense. Moreover, if \(\mu\) is a submersion at \(0\), then \(0\) is a weakly regular value (see \parencite[Th. 3.5.4]{alma994602424702061}), but not every weakly regular value is also regular. It remains to somehow motivate the slice condition. That is precisely what will be the substitute of the quotient manifold theorem. Indeed, \parencite[Sec. 3.5]{alma994602424702061} gives a detailed exposition of the theory of quotient Banach-manifolds. Under the hypothesis of the existence of a slice, we may use it to get local coordinates around every point (see the proof of \parencite[Th. 3.5.25 and Lemma 3.5.26]{alma994602424702061}). \\

Moreover, observe that if \(\phi\colon G \times \mathcal{M} \to \mathcal{M}\) is the action, then \(\phi^{(x)}\colon G \to \mathcal{M}\) is an embedding for each \(x \in \mathcal{M}\), and if \(x \in \mu^{-1}(0)\), then \(G\cdot x \subset \mu^{-1}(0)\) is an embedded submanifold. If the slice condition is satisfied, we have \( T_x(\mu^{-1}(0))=T_xS_x\oplus T_x (G\cdot x)\), hence \(T_x(\mu^{-1}(0))/T_x(G \cdot x)\cong T_xS_x\). Moreover, for each sufficiently small slice \parencite[h. 3.5.25 and Lemma 3.5.26]{alma994602424702061} state that the map \(\varphi\colon S_x \to \mathcal{M}_{\mathrm{red}}, p \mapsto \pi(p)\) is a diffeomorphism onto its image, therefore \(T_xS_x \cong T_{\pi(x)} \mathcal{M}_{\mathrm{red}}\). Putting everything together, we get:
\begin{equation} \label{eq_tangent_to_reduced_M}
   T_{\pi(x)} \mathcal{M}_{\mathrm{red}} \cong T_x \ \mu^{-1}(0)/T_x(G \cdot x) 
\end{equation}
Where the isomorphism is given by the first isomorphism theorem. Indeed, \(\ker \diff_x \pi=T_x(G\cdot x)\), hence we get an isomorphism \(\overline{\diff_x \pi}\colon T_x \ \mu^{-1}(0)/T_x(G \cdot x)\to T_{\pi(x)}\mathcal{M}_{\mathrm{red}}\). \\

Observe as well that the fact that \(\phi^{(x)}\) is an embedding means that \(T_{x}(G \cdot x)=\diff_e\phi^{(x)}(\mathrm{Lie}(G))\), thus, by equation \eqref{def_hat_X}
\begin{equation} \label{eq_tang_orbit}
T_x(G \cdot x)=\corch{X_{\xi}: \xi \in \mathrm{Lie}(G)}.
\end{equation}

\begin{obs}
    Recall that the Hausdorff condition is something imposed to manifolds to avoid pathological situations, but it is not implied by the condition of being locally euclidean. When the action is proper, we have a topological group acting continuously and properly on a topological space. It is a standard result (see for example \parencite[Prop. 21.4]{lee_introSM}) that under those assumptions the quotient space is Hausdorff.
\end{obs}

We can now state the reduction theorem for Banach manifolds.

\Theorem{\textbf{Banach symplectic reduction}}{Let \(\mathcal{M}\) be a symplectic manifold endowed with a symplectic form \(\Omega\), and a smooth, symplectic, left action of a Banach-Lie group \(G\). Suppose that \(\mu\colon \mathcal{M} \to \mathrm{Lie}(G)^*\) is a moment map for which \(0\) is a weak regular value, and satisfying the slice condition. Then, \(\mathcal{M}_{\mathrm{red}}:=\mu^{-1}(0)/G\) has the structure of a symplectic (possibly non-Hausdorff) Banach manifold, which is endowed with a symplectic form \(\Omega_{\mathrm{red}}\) and satisfying
\[
\pi^* \Omega_{\mathrm{red}}=i^*\Omega
\]
where \(i \colon \mu^{-1}(0) \hookrightarrow \mathcal{M}\) is the inclusion and \(\pi\colon \mu^{-1}(0) \to \mathcal{M}_{\mathrm{red}}\) the natural projection. Moreover, if the restricted action of \(G\) on \(\mu^{-1}(0)\) is proper, then \(\mathcal{M}_{\mathrm{red}}\) is Hausdorff.  \label{th_banach_sympl_red}
}

Before giving the proof, let us state and show some lemmas we will need. Recall that, given a skew-symmetric bilinear map \(\Omega\colon E \times E \to \R\) such that \(\Omega(u,v)=0\) for all \(v \in E\) implies \(u=0\), (that is, \(\ker \Omega^\flat=\corch{0}\)) we say that it is non-degenerate. 

\Lemma{}{Let \(E\) be a Banach space and let \(\Omega\colon E \times E \to \R\) be a skew-symmetric, non-degenerate, continuous bilinear map. Consider, for each \(v \in E\) the maps \(p_v \colon E \to \R, u \mapsto |\Omega(u,v)|\). Then, \(\corch{p_v}_{v \in E}\) is a family of seminorms with respect to which \(E\) becomes a Hausdorff locally convex TVS. Moreover, given a linear map \(\alpha\colon E \to\R\) continuous with respect to the locally convex topology, there exists \(v \in E\) such that \(\alpha(u)=\Omega(u,v)\) for all \(u \in E\). \label{lemma_banach_sympl_red}}
\begin{proof}
    The fact that \(\corch{p_v}_{v \in E}\) is a family of seminorms on \(E\) follows immediately from the fact that \(\Omega\) is bilinear. Moreover, non-degeneracy of \(\Omega\) implies that if \(p_v(u)=0\) for all \(v\in E\), then \(u=0\). Endow \(E\) with the initial topology with respect to the natural projections \(q_i\colon E\to E/p_i^{-1}(0)\). This turns \(E\) into a Hausdorff TVS. \\

    For the second statement, let \(\alpha\colon E \to \R\) be continuous with respect to the TVS topology we just defined. It is a general fact of locally convex vector spaces that a linear map \(\alpha \colon E \to \R\) is continuous if and only if there are seminorms \(\corch{p_{i_1}, \ldots p_{i_n}}\) and a constant \(C>0\) such that \(|\alpha(v)| \leq C \cdot \max_{1 \leq j \leq n} \corch{p_{i_j}(v)}\) (see \parencite[Lemma 1.4.2 and Lemma 3.4.2]{Andreas}). Back to our setting, this means that there exist \(v_1, \ldots v_n \in E\) and a constant \(C>0\) such that 
    \[
    |\alpha(v)| \leq C \cdot  \max_{1 \leq j \leq n} \corch{p_{v_j}(v)}=C \cdot \max_{1 \leq j \leq n} \corch{|\Omega(v_j, v)|}.
    \]
    Observe that \(\alpha\) vanishes on \(F=\bigcap_{i=1}^n\ker \Omega^\flat(v_j)=(\mathrm{span}(v_1, \ldots v_n))^\Omega.\) Moreover, since each \(\Omega^\flat(v_j)\) is continuous, \(F\) is closed with respect to the norm-topology and its codimension (the dimension of the quotient \(E/F\)) is at most \(n\). Choose any algebraic complement \(W\) of \(F\) in \(E\). That is, as vector spaces, \(W \oplus F=E\). Since \(W \cong E/F\), it must be finite dimensional, hence closed in \(E\) with respect to the Banach topology. This means that \(W^*\) is also finite-dimensional, and spanned by \(\Omega^\flat(v_1)|_W, \ldots \Omega^\flat(v_n)|_W\). Then \(\alpha |_W \in W^*\) can be written as \(\alpha|_W=\sum_{i=1}^n \lambda_n \cdot \Omega^\flat(v_n)|_W\) for some real coefficients \(\lambda_n\). Since in \(F\) both sides are identically \(0\), this linear combination can be extended to all of \(E\), thus
    \[
    \alpha = \sum_{i=1}^n \lambda_i \cdot \Omega^\flat(v_i)= \Omega^\flat\parent{\sum_{i=1}^n \lambda_i \cdot v_i}=\Omega^\flat(v),
    \]
    where we defined \(v=\sum_{i=1}^n \lambda_i \cdot v_i\). This proves the result.
\end{proof}

\Lemma{}{Let \(E\) be a Banach space, and let \(\Omega\colon E \times E \to \R\) be a skew-symmetric, non-degenerate, continuous bilinear map. For any closed subspace \(Y \subset E\), define the \(\Omega\)\textit{orthogonal complement}\index{\(\Omega\)-orthogonal complement} of \(Y\) by
\[
Y^\Omega:=\corch{v \in E: \Omega(u,v)=0 \ \ \text{for all} \ \ v \in Y}.
\]
Then, \(\parent{Y^\Omega}^\Omega=Y\). \label{lemma_Y_Omega_Omega}
}
\begin{proof}
    The inclusion \(Y \subset (Y^\Omega)^\Omega\) follows directly from the definition. For the reverse inclusion, let \(v \notin Y\). By a standard consequence of the Hanh-Banach theorem for locally convex TVS, there exists \(\alpha \colon E \to \R\) (continuous with respect to the locally convex topology) such that \(\alpha|_Y\equiv0\) and \(\alpha(v) \neq 0\). By the previous lemma, there is an element \(z \in E\) such that \(\alpha(u)=\Omega(u,z)\). This means that \(\Omega(u,v)=\alpha(v) \neq 0\). However, \(z \in Y^\Omega\) because for every \(y \in Y\), \(0=\alpha(y)=\Omega(y,z)=0\). Therefore, \(\alpha(v)=\Omega(v,z) \neq 0\), so \(v \notin (Y^\Omega)^\Omega\).
\end{proof}

\Lemma{}{Let \(\mathcal{M}\) and \(\mathcal{N}\) be Banach manifolds, and let \(f\colon \mathcal{M} \to \mathcal{N}\) be a surjective smooth submersion. Let \(\alpha \in \Omega^k_\R(N)\). If \(f^*\alpha=0\), then \(\alpha=0\). \label{lemma_null_pullback}
}
\begin{proof}
    \(\alpha=0\) if and only if it vanishes when applied to any set of tangent vectors at any point. Let \(p \in N\) and \(q \in M\) such that \(f(q)=p\). Choose any vectors \(v_1, \ldots v_k \in T_qN\), and \(w_1, \ldots w_k \in T_pM\) such that \(\diff_qf(w_j)=v_j\). Then:
    \[
    \alpha_p(v_1, \ldots v_j)=\alpha_{f(q)}(\diff_pf(w_1), \ldots \diff_pf(w_k))=(f^*\alpha)_q(w_1, \ldots w_k)=0.
    \]
\end{proof}

\Lemma{}{Let \(\mu\) be a moment map on a Banach manifold \(\mathcal{M}\) for the left smooth action of a Banach-Lie group \(G\). Then, for each \(p \in \mu^{-1}(0)\), we have \((T_p\ \mu^{-1}(0))^{\Omega_p}=T_p(G \cdot p)\) and \(T_p \ \mu^{-1}(0)=\parent{T_p (G \cdot p)}^{\Omega_p}\). \label{lemma_Omegap_ortho}}
\begin{proof}
    By lemma \ref{lemma_Y_Omega_Omega} it suffices to check the second identity. Let \(u \in T_p \mathcal{M}\) by such that \(\Omega_p(\hat{X}_p, u)=0\) \(\hat{X}_p \in T_p(G  \cdot p)\) (see \eqref{eq_tang_orbit}), and let us see that \(u\in T_p \ \mu^{-1}(0)\). Observe that
    \[
    \Omega(\hat{X}_p, u)=\parent{i_{\hat{X}}\Omega}_p(u)=\diff_p \mu^X(u)=\langle \diff _p \mu (v), X \rangle
    \]
    the latter being zero for all \(X \in \mathrm{Lie}(G)\) if and only if \(u\in \ker \diff_p\mu=T_p \ \mu^{-1}(0)\).
\end{proof}

\begin{proof}[\textit{Proof of theorem} \ref{th_banach_sympl_red}]
    We are now ready to prove the theorem of symplectic reduction for Banach manifolds. First, observe that the fact that \(\mathcal{M}_{\mathrm{red}}\) is a Banach manifold (Hausdorff if \(G\) acts properly on \(\mu^{-1}(0)\)) is taken care of by the slice condition, so we just need to concern ourselves with obtaining the symplectic form.\\

    The requirement that \(\pi\) be a smooth surjective submersion defines \(\Omega_{\mathrm{red}}\) uniquely. Firstly, there is an isomorphism \(T_{\pi(p)}\mathcal{M}_{\mathrm{red}} \cong T_p (\mu^{-1}(0))/T_p(G \cdot p)\) for all \(p \in \mu^{-1}(0)\) by \eqref{eq_tangent_to_reduced_M} let us denote equivalence classes using brackets \([\ \cdot \ ]\). Define, pointwise
    \[
    \parent{\Omega_{\mathrm{red}}}_{\pi(p)}([u], [v]):=\Omega_p(u,v) \hspace{1 cm} u,v \in T_p\ \mu^{-1}(0).
    \]
    Let us check that the previous definition does not depend on the representatives chosen. To that extent, let \([v]=[v']\) and \([u]=[u']\), where all vectors are elements of \(T_p(\mu^{-1}(0))\). Explicitly, by \eqref{eq_tang_orbit}, this means that there exist \(X, Y \in \mathrm{Lie}(G)\) such that \(u'=u+\hat{X}_p\) and \(v'=v+\hat{Y}_p\), where \(\hat{X}\) and \(\hat{Y}\) are the fundamental vector fields associated to \(X\) and \(Y\in \mathrm{Lie}(G)\). Then
    \begin{align*}
        \Omega_p(u', v') & = \Omega_p(u+\hat{X_p}, v+\hat{Y}_p) = \Omega_p(u,v) + \Omega_p(u, \hat{Y_p}) + \Omega(\hat{X}_p, v) + \Omega(\hat{X}_p, \hat{Y}_p) \\
        & = \Omega_p(u,v),
    \end{align*}
    where we used that \(\Omega_p(\hat{X}_p, \hat{Y}_p) = \Omega_p(u, \hat{Y}_p)=\Omega_p(\hat{X}_p, v)=0\) because, by lemma \ref{lemma_Omegap_ortho} \(T_p(\mu^{-1}(0))\) is the \(\Omega_p\)-orthogonal complement of \(T_p(G \cdot p)\). Observe that \(T_p(G \cdot p) \subset T_p \ \mu^{-1}(0)=\parent{T_p (G \cdot p)}^{\Omega_p}\) so elements of \(T_p(G \cdot p)\) are \(\Omega_p\)-orthogonal to themselves. Let us now see that the definition is also independent of the point \(p\) chosen to define \(\Omega_p\). Suppose that \(\pi(p)=\pi(p')\). Since \(\pi\) is the natural projection induced by the action \(\phi\colon G \times M \to M\), this means that there is some \(g \in G\) such that \(\phi_g(p)=p'\). Now, since \(\pi(p)=\pi(p')\), we obviously have \(T_{\pi(p)}\mathcal{M}_{\mathrm{red}}=T_{\pi(p')}\mathcal{M}_{\mathrm{red}}\), but then, the isomorphism \eqref{eq_tangent_to_reduced_M} applies for both \(p\) and \(p'\). We get that given \(u \in T_p \ \mu^{-1}(0)\) and \(u'\in T_{p'} \ \mu^{-1}(0)\), \([u]\) and \([u']\) correspond to the same element in \(T_{\pi(p)}\mathcal{M}_{\mathrm{red}}\) if and only if \(\diff_p\pi(u)=\diff_p \pi(u')\) where \(p'=\phi_g(p)\) (see the comments bellow equation \eqref{eq_tangent_to_reduced_M}). 
    Then, since the action is symplectic
    \begin{align*}
        (\Omega_{\mathrm{red}})_{\pi(p')}\big([\diff_p \phi_g(u)], [\diff_p \phi_g(v)]\big) & = \Omega_{p'}(\diff_p \phi_g(u),\diff_p\phi_g(v))=\Omega_{\phi_g(p)}(\diff_p \phi (u), \diff_p\phi(v))=\parent{\phi_g^*\Omega}_p(u,v) \\ 
        &= \Omega_p(u,v)  = (\Omega_{\mathrm{red}})_{\pi(p)}([u], [v]).
    \end{align*}
    Where we used the fact that \(\diff_{p'} \pi\parent{\diff_p\phi_g(u)}=\diff_p(\pi \circ \phi_g)(u)=\diff_p\pi(u)\) to deduce that \([u]\) and \([\diff_p\phi_g(u)]\) represent the same element in \(T_{\pi(p)}\mathcal{M}_{\mathrm{red}}\). \\
    
    This means that \(\Omega_{\mathrm{red}}\) is well-defined. It is smooth because the definition implies that \(\pi^*\Omega_{\mathrm{red}}=i^*\Omega\), and the right-hand side is smooth. Moreover, it is closed because \(\diff\) commutes with pullbacks
    \[
    \pi^*(\diff \Omega_{\mathrm{red}})=\diff\parent{\pi^*\Omega_\mathrm{red}}=\diff\parent{i^*\Omega}=i^*\parent{\diff \Omega}=0,
    \]
    and \(\pi\) is a surjective smooth submersion, so we may apply lemma \ref{lemma_null_pullback}. It remains to show that \(\Omega\) is non-degenerate. This is a linear issue, so choose \(p \in \mu^{-1}(0)\), and suppose that there is some \([u]\in T_{\pi(p)} \mathcal{M}_{\mathrm{red}}\) such that for every \([v]\in T_{\pi(p)}\mathcal{M}\) we have \((\Omega_{\mathrm{red}})_{\pi(p)}([u], [v])=0\), let us see that \([u]=[0]\). The definition of \(\Omega_{\mathrm{red}}\) allows us to translate this problem to one on \(T_p(\mu^{-1}(0))\). That is, given \(u \in T_p(\mu^{-1}(0))\) such that for every \(v \in T_p(\mu^{-1}(0))\) we have \(\Omega_p(u,v)=0\), we have to show that \(u \in T_p(G \cdot p)\). Let us apply lemma \ref{lemma_banach_sympl_red}. Take \(T_pM\) as the Banach space (it is a Banach space by lemma \ref{Lemma_tang_space_to_Banach_manifold}), and take \(T_p(G \cdot p)\) as the closed subspace. Now, lemma \ref{lemma_banach_sympl_red} implies that 
    \[
    T_p(G\cdot p)=\parent{(T_p(G \cdot p))^\Omega}^\Omega=(T_p \ \mu^{-1}(0))^\Omega=\corch{u \in T_pM: \Omega_p(u,v)=0 \ \ \text{ for all }\ \ v \in T_p \ \mu^{-1}(0)},
    \]
    which is what we wanted to show.
\end{proof}
\section{The gauge group of a principal bundle as a Hilbert Lie group} \label{section_gauge_group_is_Hilbert}
The aim of this section is to show how the gauge group of a principal bundle can be seen as a Banach manifold. To so, we shall make use of the isomorphism \(\mathscr{G}\cong \Gamma(M; \mathrm{Ad}(P))\) given by theorem \ref{equivalences_gauge_group}. It is the completion of \(\Gamma(M; \mathrm{Ad}(P))\) with respect to a Sobolev norm what actually \textit{is} a Hilbert Lie group.
\subsection{Sobolev spaces} \label{section_Sobolev_spaces} {\color{violet}
Let us consider a vector bundle \(E\), endowed with a positive-definite inner product \(\langle \cdot, \cdot \rangle\), which induces a norm. Fix a connection \(\nabla^E\) on \(E\), that is compatible with the fiber metric. This connection induces a covariant derivative in all tensor bundles over \(M\) as we show bellow (we follow the exposition in \parencite{wehrheim2004uhlenbeck}). \\

As we know, the connection \(\nabla^E\), defines a map \(\nabla^E\colon \Gamma(E) \to \Gamma(T^*M \otimes E)\). On the other hand, since \((M,g)\) is a riemannian metric, \(TM\) is naturally endowed with the Levi-Civita connection \(\nabla^g\). This induces a connection on the dual bundle \(\nabla^{g*}\). The bundle \(T^*M \otimes E\) is therefore endowed with the connection \(\nabla^{g*} \otimes \nabla^E \colon \Gamma(T^*M \otimes E) \to \Gamma(T^*M \otimes T^*M \otimes E)\). Now, \(\otimes^2 T^*M \otimes E\) is endowed with a connection \(\nabla^{g*} \otimes \nabla^E\), and we may use it to define a connection on \(T^*M \otimes \parent{\otimes^2 T^*M} \otimes E \cong\otimes^3 T^*M \otimes E\), namely, we define the connection \(\nabla^{g*} \otimes \parent{\nabla^{g*} \otimes \nabla^E}\) (see section \ref{sec_induced_connections}). Iterating this process, we may define a connection on each \(\otimes^k T^*M\otimes E\) by 
\[
\nabla^{g*} \otimes \overset{k}{\cdots} \otimes \nabla^{g*} \otimes \nabla^E.
\]
Making an abuse of notation, we will denote by \(\nabla\) all the so-defined connections (including \(\nabla^E\)). Explicitly, if \(\alpha \in \Gamma(\otimes^k T^*M \otimes E)\), then \(\nabla \alpha \in \Gamma(\otimes^{k+1}T^*M \otimes E)\) is defined by \parencite[Eq. B2]{wehrheim2004uhlenbeck}
\begin{align*}
    \nabla \alpha(X_0, \ldots X_k)& =  \nabla^E(\alpha(X_1, \ldots X_k)) - \alpha(\nabla_{X_0}^g X_1, \ldots X_k) - \ldots - \alpha(X_1, \ldots \nabla_{X_0}^g X_k), 
\end{align*}

For \(k \in \N_{\geq 1}\) and any two sections \(\phi, \psi \in \Gamma(E)\), we define the \(L_k^2\) inner product
\begin{equation} \label{prod_in_Gamma(E)}
\langle \phi, \psi\rangle_k:= \int_M \bigg( \langle \phi, \psi \rangle + \langle \nabla \phi, \nabla \psi \rangle  + \ldots+ \langle \nabla^k \phi, \nabla^k \psi \rangle \bigg) \ \diff V_g.
\end{equation}

\begin{obs}
    Each term in the sum inside the integral is a smooth function, so it makes sense to integrate it with respect to the volume form \(\diff V_g\), and the compactness of \(M\) ensures that the integral will be finite. As for the first claim, the symbol \(\langle \nabla^\ell \phi, \nabla^\ell\psi \rangle\) \textit{hides} the fact that we need to define an inner product on tensor products. First, the fact that \(g\) is non-degenerate (it is a riemannian metric) allows us to define an inner product of one-forms by \(\langle \alpha, \beta\rangle_g:=g\parent{\alpha^\#, \beta^\#}\), where \(\#\) is the \textit{sharp} isomorphism (see \parencite[Ch. 2]{lee_introRM}). On the other hand, whenever \(V_1, \ldots V_n\) are vector spaces endowed with inner products \(\langle\cdot, \cdot \rangle_i\), we may define an inner product by extending by linearity the following expression
    \[
    \langle v_1 \otimes v_2 \otimes \cdots \otimes v_n,  w_1 \otimes w_2 \otimes \cdots \otimes w_n \rangle = \langle v_1, w_1 \rangle_1 \cdot \langle v_2, w_2 \rangle_2 \cdots \langle v_n, w_n \rangle_n.
    \]
    This defines an inner product on each \(\otimes^k T^*M \otimes E\) (see \parencite[Lemma 1.7]{Rusman} for a very detailed proof that the previous expression defines an inner product).
\end{obs}
The map
\[
\begin{array}{ccc}
    \Gamma(E) \times \Gamma(E) & \longrightarrow & \R \\
    (\phi, \psi) & \mapsto & \langle \phi, \psi\rangle_k
\end{array}
\]
defines an inner product on \(\Gamma(E)\). Of course, everything works out fine if we work with complex vector bundles and a hermitian fiber metric, in which case we can define a hermitian inner product on \(\Gamma(E)\). Since \(\Gamma(E)\) is a vector space with an inner product, we can always consider its Hilbert completion, which we shall denote \(L_k^2(M;E)\) or \(L_2^k(E)\). These spaces are called \textit{Sobolev spaces}\index{Sobolev space}.
}

% \Lemma{}{Let \((M,g)\) be a compact riemannian manifold, and let \(E\) be a vector bundle over \(M\) endowed with a fiber product \(\langle \cdot, \cdot \rangle\). The map defined in equation \ref{prod_in_Gamma(E)}: 
% \[
% \begin{array}{ccc}
%     \Gamma(E) \times \Gamma(E) & \longrightarrow & \R \\
%     (\phi, \psi) & \mapsto & \langle \phi, \psi\rangle_k
% \end{array}
% \]
% is an inner product in the vector space \(\Gamma(E)\).
% }
% \begin{proof}
%     It follows straightforward from the fact that the product \(\langle \cdot, \cdot \rangle\) on each summand is pointwise an inner product and the fact that the integral of a non-negative function with respect to the riemannian volume form is zero if and only if the function itself is zero \parencite[Prop. 16.28]{lee_introSM}.
% \end{proof}
% Observe that we do not need the trick we use in \(L^p\)-spaces of identifying functions that coincide almost everywhere to ensure positivity because we do not have the problem of section being zero almost everywhere. Indeed, sections are smooth, and if they are zero on a dense subset, then they are identically zero. \\

% With the previous constructions in mind, we can give the final definition of the Sobolev space associated to the space of connections.

% \Definition{}{Let \((M, g)\) be a compact riemannian manifold, and let \(E\) be a vector bundle over \(M\) endowed with an inner product and a covariant derivative. Let \(m \in \N_{\geq 0}\) and \(1 \leq p <\infty\). The \textit{Sobolev space} \(W^k(M; E)\) is the Hilbert-completion of \(\Gamma(M;E)\) with respect to the inner product \(\langle \cdot, \cdot \rangle_k\).}

% Of course, once the completion is then endowed with the \(\langle \cdot, \cdot \rangle_k\), and this turns \(W^{k}(M;E)\) into a Hilbert space. 

\begin{obs}
    Recall that the topology of a Hilbert space is the metric topology with respect to the distance induced by the norm. The inner product we defined depends on \(g, \nabla\) and \(\langle\cdot, \cdot \rangle\). However, different choices of these parameters lead to equivalent norms. This means that, as a topological space, \(L^2_k(M;E)\) is well-defined and independent of any choice of these parameters.
    % \begin{enumerate}
    %     \item The previous construction can be done using Banach spaces. There, instead of an inner product, we define a norm
    %     \[
    %     ||\phi||_{p,k}:=\parent{\int_M \parent{||\phi||_p+||\nabla\phi||_p + \ldots ||\nabla^k \phi ||_p} \ \diff V_g}^{1/p}
    %     \]
    %     out of a \(p\)-norm on each \(\otimes^k T^*M \otimes E\), and then we consider \(W^{p,k}(M;E)\) as the completion of \(\Gamma(E)\) with respect to that norm.
    %     \item We obviously have \(||\phi||_k < ||\phi||_{k'}\) for every \(k<k'\) (the latter has more summands).
    %     \item Recall that the topology of a Hilbert space is the metric topology with respect to the distance induced by the norm. The inner product we defined depends on \(g, \nabla\) and \(\langle\cdot, \cdot \rangle\). However, different choices of these parameters lead to equivalent norms. This means that, as a topological space, \(W^k(M;E)\) is well-defined and independent of any choice of these parameters.
    %     % \item There is a way to write the norm in a more explicit way. Let \((U_i, \eta_i)_{i=1}^m\) be a finite covering of \(M\) by local charts (which exists by compactness of \(M\)), and such that there are trivializations \(\Phi_{U_i}\colon \pi^{-1}(U_i) \to U_i \times \R^k\) (or \(\C^k\) is the bundle is complex). Then, the product \(\langle \cdot, \cdot \rangle_k\) is equivalent to the product defined by
    %     % \[
    %     % \langle \phi, \psi \rangle'_k:=\sum_{i=1}^m \sum_{|\nu| \leq k} \int_{\eta_i^{-1}(U_i)} \bigg\langle \frac{\partial^{|\nu|}(\mathrm{pr}_{\R^k} \circ \Phi_{U_i} \circ \phi \circ \eta_i)}{\partial x^\nu }, \frac{\partial^{|\nu|}(\mathrm{pr}_{\R^k} \circ \Phi_{U_i} \circ\psi \circ \eta_i)}{\partial x^\nu } \bigg\rangle \ \diff ^n \bm x.
    %     % \]
    %     % where we are considering derivatives of real (or complex) valued functions, and the inner product is the usual euclidean (or hermitian) inner product (see \parencite[Section 5.7]{Rudolph_2017pug}).
    % \end{enumerate}
\end{obs}
Elements of a Hilbert completions are, formally, equivalence classes of Cauchy sequences. The following lemma provides a more tangible way of looking at these spaces.

\Lemma{\textbf{Sobolev}}{Let \(E \to M\) be a vector bundle over a compact smooth manifold, and let \(k,p \in \N_{\geq 0}\) be such that \(k> \dim M/2+p\). Then, there is a continuous embedding \(W^k(M;E) \hookrightarrow \mathscr{C}^p(M;E)\). \label{Sobolev_lemma} }
\begin{proof}
    See \parencite[Th. B.2]{wehrheim2004uhlenbeck} Here, \(\mathscr{C}^p(M;E)\) is endowed with the metric topology induced by the norm 
    \[
    ||\phi||_\infty:=\sup_{j \leq p} \big\{ \max_{p \in M}\ \langle \nabla^j \phi(p), \nabla^j \phi(p)\rangle \big\}. 
    \]
\end{proof}
That is, Sobolev spaces of sections can be seen as spaces of usual sections but with the caveat that some regularity may be lost. \\

We present now a sketch of the construction that allows us to see \(\mathscr{G}\) as a Hilbert Lie group. For the construction, we consider \(G\) a compact, connected, semisimple \textit{matrix Lie group}. That is, \(G \subset \mathfrak{gl}(n,\C)=\mathrm{M}_n(\C)\) for some \(n\in \N\). We suppose that \(M\) is a compact smooth oriented riemannian manifold. Consider the vector bundle \(E^n_\C:=P \times_G \mathfrak{gl}(n, \C)\) where \(G\) acts on \(\mathfrak{gl}(n, \C)\) by conjugation. We may see \(\mathrm{Ad}(P)\) as a subbundle of \(E^n_\C\), i.e. it is an embedded submanifold, and its projection is the restriction of the projection of \(E_\C^n\). Observe that \(\mathrm{Ad}(P)\) is also constructed using the action of \(G\) by conjugation. Using the riemannian metric and any hermitian fiber metric on \(E_\C^n\) (those always exist, and they can be constructed through a partition of unity argument, see \parencite[Page 196]{lee_introCM}). Hence, we can consider \(\Gamma(M; \mathrm{Ad}(P))\subset \Gamma(M; E^n_\C)\subset L_k^2(M; E_\C^n)\). We define \(\mathscr{G}_k:=\overline{\Gamma(M;\mathrm{Ad}(P))}\), i.e. the closure of \(\Gamma(M;\mathrm{Ad}(P))\) in \(L_k^2(M;E_\C^n)\) with respect to the Sobolev topology of the latter. If we choose \(k \in \N_{\geq 0}\) satisfying \(k>\dim M/2+1\), then elements of \(\mathscr{G}_k\) can be seen as \(\mathscr{C}^2\) sections of the bundle \(\mathrm{Ad}(P)\). Hence, \(\mathscr{G}_k\) has a natural group structure where the group operation is defined pointwise. Actually, \(L_k^2(M; E_\C^n)\) has a natural product structure. This product, call it \(m\colon L_k^2(M; E_\C^n) \times L_k^2(M; E_\C^n) \to L_k^2(M; E_\C^n) \) is bilinear and continuous (see \parencite[Eq. 2.2]{mitter1981bundle}). Hence, \(m(\overline{\Gamma(M; \mathrm{Ad}(P))\times \Gamma(M; \mathrm{Ad}(P)}))\subset \overline{\Gamma(M; \mathrm{Ad}(P))}\), and \(\mathscr{G}_k\) is closed under products. 



% As a consequence, even though \(W^k(M;E)\) is formally constructed using equivalence classes of Cauchy sequences, we can still view its elements as maps (sections in fact) \(M \to E\), but now some of those elements lose regularity; we can only ensure they are of class \(\mathscr{C}^p\). \\

% We are now ready to endow \(\mathcal{G}\) with a Hilbert-group structure. To keep the technicalities to a minimum, we will suppose that \(G\) is a compact, connected, matrix Lie group, and that \(M\) is a compact manifold. In that case, there is \(n \in \N_{\geq 1}\) such that \(G \subset \mathfrak{gl}(n, \C)=\mathrm{M}_n(\C)\). Consider the vector bundle \(E^n_\C:=P \times_G \mathfrak{gl}(n, \C)\) where \(G\) acts on \(\mathfrak{gl}(n, \C)\) by conjugation. Consider as well the bundle \(\mathrm{Ad}(P)=P \times_G G\) where \(G\) acts on itself by conjugation again. \(\mathrm{Ad}(P)\) is a subbundle of \(E_\C^n\), when both of them are seen as and therefore, we have \(\Gamma(M; \mathrm{Ad}(P)) \subset \Gamma(M; E_\C^n)\). Now, \(M\) always admits a riemannian metric (see \parencite[Prop. 13.3]{lee_introSM}), and  every vector bundle admits a fibre product and a connection compatible with it (this is a partition of unity argument, see for example \parencite[Page 196]{lee_introCM}). Therefore, we may consider \(\Gamma(M; E_{\C}^n) \subset W^k(M; E_\C^n)\) for any \(k > \dim M/2+1\). In particular, we can see \(\Gamma(M; \mathrm{Ad}(P))\) as a subset of the Hilbert space \(W^{k+1}(M; E_\C^n)\). We can now define \(\mathscr{G}_k:=\overline{\Gamma(M; \mathrm{Ad}(P))}\), where the closure is taken with respect to the metric topology in \(W^{k+1}(M, E_\C^n)\). 

% \Lemma{}{With the previous definitions, \(\mathrm{Ad}(P)\) and \(\mathrm{ad}(P)\) are subbundles of \(E_\C^n\).}
% \begin{proof}
%     Since \(G\) is a matrix Lie group, we may see it as an embedded submanifold of \(\mathfrak{gl}(n, \C)=\mathrm{M}_n(\C) \cong \C^{n^2}\). Moreover, in that case, its lie algebra is a matrix lie algebra, which consists of square matrices of the same dimension. Therefore, \(\mathfrak{g}\) is a vector subspace of \
%     \(\mathfrak{gl}(n, \C)\), and therefore an embedded submanifold. Moreover, all the bundles \(\mathrm{Ad}(P), \  \mathrm{ad}(P)\) and \(E_\C^n\) are associated bundles obtained from, \(P \times G\), \(P \times \mathfrak{g}\) and \(P \times \mathfrak{gl}(n, \C)\) respectively. To get the associated bundle, \(G\) acts by conjugation on \(G\) and \(\mathfrak{gl}(n, \C)\), and through the \(\mathrm{Ad}\)-representation on \(\mathfrak{g}\). However, since we work with matrix Lie groups, the adjoint representation is just conjugation, so all associated bundles are obtained from the action of \(G\) on the corresponding space by conjugation. Finally, seeing \(G\) and \(\mathfrak{g}\) inside \(\mathfrak{gl}(n, \C)\), the equivalence relation defining the associated bundles does not add elements from to the orbits from \(\mathfrak{gl}(n, \C)\) that were not present before (i.e. it cannot happen that for \(h \notin \mathfrak{g}\), we end up having \([p,h] \in \mathrm{ad}(P)\)). Thus, we can see \(\mathrm{Ad}(P)\) and \(\mathrm{ad}(P)\) as {\color{red}embedded submanifolds} of \(E_\C^n\) through the inclusion \(\imath\colon P \times_G G \hookrightarrow P \times _G \mathfrak{gl}(n, \C), [p, g] \mapsto [p,g]\), the latter being an equivalence class in \(E_\C^n\). In particular, \(\Gamma(M; \mathrm{Ad}(P)) \subset \Gamma(M; E_\C^n)\).
% \end{proof}

% \Lemma{}{With the previous definitions, choose \(k \in \N_{\geq 0}\) such that \(k>\dim M/2+1\). Then, \(\mathscr{G}_k\) can be seen as a space of \(\mathscr{C}^2\)-sections of the bundle \(\mathrm{Ad}(P)\), which we will denote \(\Gamma^2(M;\mathrm{Ad(P)})\), and it is a group under pointwise multiplication. \hfill \({\color{blue} \bigstar}\)}
% \begin{proof}
%     Let us see first that maps in \(\mathscr{G}_k\) are indeed sections. More precisely, let \(\phi \in \mathscr{G}_k\). Lemma \ref{Sobolev_lemma} ensures that 
%     \(\phi\) can be seen as a \(\mathscr{C}^2\)-map from \(M\) to \(E_\C^n\). By definition of the Hilbert-completion, and due to the fact that the inclusion given by lemma \ref{Sobolev_lemma} is continuous, there is a sequence \((\phi_n)_{n \in \N} \in \Gamma(M;\mathrm{Ad}(P))\) satisfying \(||\phi_n-\phi||_\infty \xrightarrow{n \to \infty}0\). By definition of the \(||\cdot ||_\infty\)-norm, norm convergence implies pointwise convergence, hence \(\phi_n(p) \xrightarrow{n \to \infty} \phi(p)\) for every \(p \in M\). Finally, let \(\pi\colon E_\C^n \to M\) be the vector bundle projection. Since it is continuous, and \(\phi_n(p)\) converges in the topology of \(E\), we have
%     \[
%     p=\lim_{n \to \infty} \pi(\phi_n(p))=\pi\parent{\lim_{n \to \infty} \phi_n(p)}=\pi(\phi(p)),
%     \]
%     and therefore, \(\pi \circ \phi = \mathrm{id}_M\), so \(\phi\) is a section of the vector bundle \(E_\C^n\). Moreover, \(\phi(p) \in \mathrm{Ad}(P)\) for every \(p \in M\), so \(\phi\) is actually a section of \(\mathrm{Ad}(P)\). {\color{red} revisar. Para esto tengo que especificar bien que significa que \(\mathrm{Ad}(P)\) sea un vertical subbundle. No es un fibrado vectorial.} \\

%     Once established that \(\mathscr{G}_k\) is the space of \(\mathscr{C}^2\)-sections, the product of two sections can be defined pointwise just as in proposition \ref{equivalences_gauge_group}. Moreover, it can be shown that for a compact manifold, if \(\phi, \psi \in W^{k+1}(M;E_\C^n)\), then its pointwise product \(\phi \cdot \psi\) does not lose regularity, i.e. \(\psi \cdot \phi \in W^{k+1}(M;E_\C^n)\) (see \parencite[Lemma B.3]{wehrheim2004uhlenbeck}). In particular, the pointwise product of two sections in \(\mathscr{G}_k\) still belongs to \(\mathscr{G}_k\), so this set is a group.  
% \end{proof}

% We are now ready to proof the desired result.

\Theorem{}{Let \(P(M,G)\) be a principal \(G\)-bundle over a compact manifold, where \(G\) is a compact, connected, matrix Lie group. Then, for \(k>\dim M/2+1\), \(\mathscr{G}_k\) is a Hilbert Lie group whose lie algebra is \(\mathrm{Lie}(\mathscr{G}_k)\cong L^2_k(M;\mathrm{ad}(P))\), and whose exponential map is given by
\[
\exp_{\mathscr{G}_k}(\xi)(p) = \big[p,\exp_G(\xi(p))\big], \hspace{1 cm} \xi \in \mathrm{Lie}(\mathscr{G}_k), \  p \in M
\]
}
\begin{proof}
    See \parencite{mitter1981bundle}.
    % We will follow the original proof found in sections I and II of \parencite{mitter1981bundle}. We have seen that \(\mathscr{G}_k\) is a group, and it is a topological space with respect to the subspace topology inherited from \(W^{k+1}(M;E_\C^n)\). Since the latter is a Hilbert space, it is Hausdorff, so \(\mathscr{G}_k\) is Hausdorff as well. Let us check first that the group operations are continuous. To do so, we can check that the pointwise operations are continuous in \(W^{k+1}(M;E_\C^n)\), because in that case, the restriction to \(\mathscr{G}_k\) will be continuous as well by the characteristic property of the subspace topology. Now, \parencite[Lemma B.5]{wehrheim2004uhlenbeck} also justifies that there is a constant \(C>0\) such that \(||\phi \cdot \psi ||_k \leq C \cdot ||\phi||_k \cdot ||\psi||_k\) and products in \(W^{k+1}(M;E_\C^n)\) are bilinear (a bilinear map defined in a Banach space is continuous if and only if it is bounded). This shows that products are continuous. To check that taking inverses is continuous {\color{red}???}  \\

    % Observe that we require \(G\) to be connected becasue the codomain of \(\exp_G\) is the connected component of the identity 

    % To endow \(\mathscr{G}_k\) with the structure of a Hilbert manifold, we are going to use the exponential map as a chart. Note that if \(\xi \in W^{k+1}(M;\mathrm{ad}(P))\), then it is a section \(M \to P\times_G \mathfrak{g}\). If, at each point we consider the unique representative such that \(\xi(p)=[p, X]\), and we identify \(X\) with \(\xi(p)\), then we can make \(\exp_G\) act on the second component, and we get an element of \(G\). Thus,
    % \[
    % \begin{array}{cccc}
    %     \exp_{\mathscr{G}_k}(\xi) \colon & M & \longrightarrow & \mathrm{Ad}(P) \\
    %     & p & \mapsto & [p, \exp_G(\xi(p))],
    % \end{array}
    % \]
    % is well-defined. We are basically using the exponential of the original Lie group to define the new exponential map pointwise, but we make the abuse of notation of simply writing \(\exp_{\mathscr{G}_k}(\xi) = \exp_G(\xi)\). Let us see that \(\exp_{\mathscr{G}_k}\) defines a homeomorphism between a sufficiently small neighborhood \(V_k(0)\) of \(0 \in W^{k+1}(M; \mathrm{ad}(P))\), and \(N_k(e)\) of the map \(p \mapsto [p, e] \in \mathscr{G}_k \subset W^{k+1}(M; E_\C^n)\), which we will call \(e\) making some abuse of notation. To do so, observe that \(\exp_{\mathscr{G}_k}(\xi)\) can be seen as an element of \(W^{k+1}(M; E_\C^n)\), which is a Hilbert space. Since \(W^{k+1}(M; \mathrm{ad}(P))\) is a Hilbert space as well, we may see the exponential map as a map
    % \[
    % \exp_{\mathscr{G}_k} \colon W^{k+1}(M; \mathrm{ad}(P)) \longrightarrow W^{k+1}(M;E_\C^n).
    % \]
    % If we consider it as above, it makes sense to speak about its regularity. Assume for a moment that \(\exp_{\mathscr{G}_k}\) is at least of class \(\mathscr{C}^1\), and that its differential defines an isomorphism from \(W^{k+1}(M; \mathrm{ad}(P))\) onto its image, which is again \(W^{k+1}(M; \mathrm{ad}(P))\). Since \(W^{k+1}(M; E_\C^n)\) is a Hilbert space, every closed subspace is complemented by its orthogonal complement. Therefore, the Banach version of the immersion theorem applies (theorem \ref{Local_immersion_Banach_Spaces}), and \(\exp\) defines a homeomorphism from \(V_k(0)\) to its image, call it \(N_k(e)\), which is an open subset of \(\mathscr{G}_k\). We thus have a chart. \\

    % Let us check now that \(\exp\) is smooth. {\color{red} Terminar.} \\

    % This allows us to proof that products are smooth near the idenitity. Indeed, if we choose \(\phi_1, \phi_2 \in V \), where \(V\) is a neighborhood of \(e\) satisfying \(V\cdot V \subset N_k(e)\), then \(V \subset N_k(e)\), and we may express \(\phi_i=\exp_{\mathscr{G}_k}(\xi_i), \ i=1,2\). Therefore, pointwise, we may use the Baker-Campbell-Hausdorff (BCH) formula, and we have
    % \[
    % \phi_1 \cdot \phi_2 = \exp_G (f(\xi_1, \xi_2)),
    % \]
    % where \(f\) is the BCH series
    % \[
    % f(\xi_1, \xi_2)=\xi_1+\xi_2+\frac{1}{2}[\xi_1, \xi_2] +\ldots
    % \]
    % where again, everything is understood pointwise. If we choose \(\xi_i\) such that \(||\xi_i||_k\) is sufficiently small, by lemma \ref{Sobolev_lemma}, we know that \(||\xi_k||_\infty\) will be small as well. In particular, we can make \(||\xi_i(p)||_\infty\) as small as we want. Now, \(G\) is a Lie group, and the exponential we are using pointwise is that of \(G\). We know from the general theory of Lie groups that the BCH series converges for sufficiently small elements in the Lie algebra. By taking \(||\xi_i||_k\) sufficiently small, we see that the series converges absolutely in \(W^{k+1}(M; \mathrm{ad}(P))\) and \(f\) is smooth. Therefore, since the smooth structure is given by the exponential map, we have
    % \[
    % \phi_1 \ \phi_2 =\exp_{\mathscr{G}_k}(f(\xi_1, \xi_2))=\exp_{\mathscr{G}_k}\big(f((\exp_{\mathscr{G}_k})^{-1}(\phi_1), (\exp_{\mathscr{G}_k})^{-1}(\phi_2))\big)
    % \]
    %  so the product is smooth near the identity.  \\

    % Keeping the previous notations, we are going to transport the smooth structure by right translation. Since we have already shown that \(\mathscr{G}_k\) is a topological group, we know that there is a neighborhood of \(e\), let us call it \(W\), satisfying \(W^{-1}\cdot W \subset N_k(e)\), and \(W \subset N_k(e)\). In particular, \(e \in W\), which means that for every \(\phi \in \mathscr{G}_k\), \(W(\phi):=W \cdot \phi\) is a neighborhood containing \(\phi\). We may define a chart on \(W\) as follows: given \(\psi\in W\cdot \phi\), there exist \(\vartheta \in W\) such that \(\psi=\vartheta \cdot \phi\). Since \(W \subset N_k(e)\), there exists \(\xi \in W^{k+1}(M;\mathrm{ad}(P))\) such that \(\vartheta = \exp_{\mathscr{G}_k}(\xi)\). Therefore, the map \(W \to V_k(0), \psi \mapsto \xi\) is a homeomorphism. Indeed, it is the composition of two homeomorphisms \((\exp_{\mathscr{G}_k})^{-1}\circ R_{\phi^{-1}}\). To check that these maps actually define charts, we need only check that the corresponding transition maps are smooth. \\

    % Let \(\psi \in W(\phi_1) \cap W(\phi_2)\), where \(W(\phi_i)=\phi_i \cdot W\). Applying the construction above, \(\psi\) is in the domain of two charts, and we may write
    % \[
    % \psi=\exp_{\mathscr{G}_k}(\xi_1)\cdot \phi_1 = \exp_{\mathscr{G}_k}(\xi_2)\cdot \phi_2.
    % \]
    % We thus have
    % \[
    % \exp_{\mathscr{G}_k}(\xi_1) = \exp_{\mathscr{G}_k}(\xi_2) \ \phi_2 \ \phi_1^{-1}.
    % \]
    % Now, observe that \(\phi_1 \ \phi_2^{-1}=\exp_{\mathscr{G}_k}(\xi_2)^{-1} \exp_{\mathscr{G}_k}(\xi_1) \in W\), so there is some \(\zeta\in V_k(0)\), which remains fixed, such that \(\phi_1 \ \phi_2^{-1}=\exp_{\mathscr{G}_k}(\zeta)\). Thus, \(\exp_{\mathscr{G}_k}(\xi_1)=\exp_{\mathscr{G}_k}\big(f(\xi_2, \zeta)\big)\), and therefore, \(\xi_1=f(\xi_2, \zeta)\), which is smooth. We have thus shown that the transition functions are smooth.

    % Moreover, the product is a smooth map. We have already shown it near the identity, and for the general case, we can repeat the argument using the transported charts. Finally, smoothness of the inversion map follows from remark \ref{Lie_group_inverse_superfluous}. \\

    % We conclude that \(\mathscr{G}_k\) is a Hilbert Lie group. Observe that since tangent spaces are isomorphic as vector spaces to the Banach spaces the manifolds are modelled on, we deduce that the lie algebra of \(\mathscr{G}_k\) is isomorphic as a Banach space to \(W^{k+1}(M;\mathrm{ad}(P))\) (see lemma \ref{Lemma_tang_space_to_Banach_manifold}). Observe that the Lie algebra structure on \(W^{k+1}(M; \mathrm{ad}(P))\) is given by the isomorphism.
\end{proof}

\begin{obs} \label{obs_def_bracket_sections}
    \(\Gamma(M; \mathrm{ad}(P))\) is a Lie algebra endowed with the pointwise Lie bracket it gets from \(\mathfrak{g}\). Given \(\bm{\xi}, \bm{\zeta}\in \Gamma(M; \mathrm{ad}(P_G))\) we define \([\bm{\xi}, \bm{\zeta}]\) as the section 
    \[
   [\bm{\xi}, \bm{\zeta}](x):= \big[\ [p, \xi], [p, \zeta] \ \big ]=[p, [\xi, \zeta]\ ] \hspace{1 cm} \text{where } \bm{\xi}(x)=[p, \xi], \bm{\zeta}(x)=[p, \zeta].
    \]
    This is well-defined independently of the chosen representatives because \(\Ad{g}\) is a Lie algebra isomorphism for all \(g \in G\).
\end{obs}

Observe that the smooth gauge group \(\Gamma(M;\mathrm{Ad}(P))\) can be seen as a subgroup of \(\mathscr{G}_k\) for each \(k\). Moreover, since the exponential map is smooth, it maps \(\Gamma(M; \mathrm{ad}(P))\) to \(\mathscr{G}\), hence the former may be seen as the Lie algebra of \(\mathscr{G}\). It is quite common in the literature to omit the Sobolev completions and directly state that \(\Gamma(M; \mathrm{ad}(P))\) is the Lie algebra of \(G\). It is useful to see \(\Gamma(M; \mathrm{ad}(P))\) from different points of view. To that end, observe that this Lie algebra is isomorphic to the following 
\[
\mathfrak{G}:=\corch{X \in \mathfrak{X}(P): (\phi_g)_*X=X \ \ \forall \ g \in G \ \ \text{and} \ \ X \  \text{ is vertical}} \hspace{1 cm}
\mathsf{G}:=\corch{\zeta \in \mathscr{C}^\infty(P; \mathfrak{g}): \zeta(p\cdot g)=\Ad{g^-1}(\zeta(p))}. 
\]
The first in endowed with the restriction of the Lie bracket it \(\mathfrak{X}(P)\), and the latter with the pointwise Lie bracket induced by \(\mathfrak{g}\). See \parencite[Sec. 3.5]{schmid2004infinite}


% \subsection{The Lie algebra of \(\mathcal{G}\)}
% {\color{red}In the previous section, we have given \(\mathscr{G}_k\) the structure of a Hilbert-Lie group. Moreover, under the idenitification \(\mathcal{G} \equiv \Gamma(M; \mathrm{Ad}(P))\), we can see \(\mathcal{G}\) as a subset (in fact, a subgroup) of \(\mathscr{G}_k\). Similarly, we can see \(\Gamma(M; \mathrm{ad}(P))\) as a subalgebra of \(\mathrm{Lie}(\mathscr{G}_k)\). In fact, if we pick \(\xi \in \Gamma(M, \mathrm{ad}(P))\), it is a smooth section, and if we consider \(\exp_{\mathscr{G}_k}(\xi)\), it is the composition of smooth maps, hence smooth, and it is an element of \(\Gamma(M; \mathrm{Ad}(P))=\mathcal{G}\). Thus, the exponential map sends \(\Gamma(M; \mathrm{ad}(P))\), so we can agree to call this algebra the Lie algebra of \(\mathcal{G}\), and denote \(\mathrm{Lie}(\mathcal{G}):=\Gamma(M; \mathrm{ad}(P))\).} \\

% Our objective in this section is to see that there is a Lie algebra isomorphism
% \[
% \mathrm{Lie}(\mathcal{G})\cong \mathfrak{G}:=\corch{X \in \mathfrak{X}(P): (\phi_g)_*X=X \ \ \forall \ g \in G, \  \text{and} \ \ X \  \text{ is vertical}  },
% \]
% where \(\phi\colon G \times P \to P\) is the right action. The latter is a Lie algebra when it is endowed with the Lie bracket in \(\mathfrak{X}(P)\) (although we must show that it is invariant under that Lie bracket). Observe that the first set is composed of vector fields \textit{defined on} \(P\), and the algebra structure we endow them with is that inherited from the Lie bracket in \(\mathfrak{X}(P)\). On the other hand, \(\Gamma(M; \mathrm{ad}(P))\) has a natural Lie algebra structure defined pointwise using the Lie bracket of \(\mathfrak{g}\). More precisely, if \(s_1, s_2 \in \Gamma(M; \mathrm{ad}(P))\), then \([s_1, s_2]\) is the section defined pointwise by
% \[
% [s_1, s_2](p):=\big[\ p, [s_1(p), s_2(p) ] \ \big].
% \]
% This is well-defined. 

% \Lemma{}{Let \(P(M,G)\) be a principal \(G\)-bundle over a manifold \(M\). Then, \((\mathfrak{G}, [\cdot, \cdot])\) is a Lie subalgebra of \((\mathfrak{X}(P), [\cdot, \cdot ]).\) \(\hfill\) \({\color{blue} \bigstar}\) 
% }
% \begin{proof}
%     Recall that vertical vector fields are those lying on the vertical bundle. Moreover, \(G\)-invariant vector fields form a vector space, so it is obvious that \(\mathfrak{G}\) is a vector space. To check that it is a subalgebra, recall that the Lie bracket of \(G\)-invariant vector fields is again \(G\)-invariant (\((\phi_g)_*[X_1, X_2]=[(\phi_g)_*X_1, (\phi_g)_*X_2]\)), so we need only check that the bracket of vertical vector fields is vertical again. Observe that \(X \in \mathfrak{X}(P)\) is vertical if and only if \(\diff_p\pi(X_p)=X_p(\pi)=0\) for every \(p \in P\), that is, \(X(\pi) \equiv 0\). Thus, let \(X,Y \in \mathfrak{X}(P)\) be vertical vector fields. Therefore:
%     \[
%     \diff_p\pi([X,Y]_p)=[X,Y]_p(\pi)=X_p(Y(\pi))-Y_p(X(\pi))=0,
%     \]
%     as we wanted to show. We thus conclude that \((\mathfrak{G}, [\cdot, \cdot])\) is a Lie algebra.
% \end{proof}

% We can finally proof the desired result.

% \Proposition{}{Let \(M\) be a smooth manifold, and \(P(M,G)\) a principal \(G\)-bundle over \(M\). Let \(\mathfrak{G}\) be the Lie algebra defined as above. Then, \((\mathfrak{G}, [\cdot, \cdot]) \cong (\mathrm{Lie}(\mathcal{G}), [\cdot, \cdot])\). \hfill \({\color{blue} \bigstar}\)}
% \begin{proof}
%     We will use an intermediate Lie algebra that is isomorphic to both of them. First, observe that every \(X \in \Gamma(\mathrm{ad}(P))\) is a map \(M \to P \times_G \mathfrak{g}, m \mapsto [p, X_p]\), where \(\pi(p)=m\), and \(X_p \in \mathfrak{g}\). By definition of the \(\mathrm{ad}(P)\)-bundle, \([p, X_p]=[p\cdot g, \Ad{g^{-1}}(X_p)]\). Thus, we can identify each section with a smooth map \(\xi\colon P \to \mathfrak{g}\) satisfying \(\xi(p\cdot g)=\Ad{g^{-1}}(\xi(p))\). The set
%     \[
%     \mathtt{G}:=\corch{\xi \in \mathscr{C}^\infty(P, \mathfrak{g}): \xi(p\cdot g)=\Ad{g^{-1}}(\xi(p))}
%     \]
%     is a Lie algebra under pointwise bracket, and it is isomorphic to \((\Gamma(\mathrm{ad}(P)), [\cdot, \cdot])\). Indeed, let \(X \in \Gamma(\mathrm{ad}(P))\), we define \(\xi(p):=X_p\), where \(X(m)=[p, X_p]\) (more precisely, \(\xi(p)=q_p^{-1}(X(m))\), see remark \ref{def_q_p}). It is clear that \(\xi(p \cdot g)=\Ad{g^{-1}}(\xi(p))\) due to the relation that defines \(\mathrm{ad}(P)\). Moreover, \(\xi\) is smooth. This can be seen by choosing a trivialization (see equation \ref{trivialization_associated_bundle} for the explicit form). We find, in the appropriate open sets, \(\xi=\Ad{s_U} \circ \mathrm{pr}_{\mathfrak{g}} \circ \Phi \circ X \circ \pi\), which is a composition of smooth maps. It helps to visualize it with a commutative diagram:
%     \begin{center}
%     \begin{tikzcd}[row sep=normal, column sep=normal, text height=1ex, text depth=1ex, ampersand replacement=\&]
%             U \subset M \arrow[r, "{\displaystyle X}"]  \& \pi_G^{-1}(U) \subset P\times \mathfrak{g} \arrow[r, "{\displaystyle \Phi}"]\& U \times \mathfrak{g} \arrow[r, "{\displaystyle \mathrm{pr}_{\mathfrak{g}}}"] \& \mathfrak{g}  \\ 
%             \\
%             \pi^{-1}(U) \subset P \arrow[uu, "{\displaystyle \pi}"] \arrow[uurrr, "{\displaystyle \Ad{s_U^{-1} \circ \xi}}", '] 
%     \end{tikzcd} 
%     \end{center}
%     Since the Lie bracket in \(\Gamma(M, \mathrm{ad}(P))\) is defined pointwise in the second coordinate, the map sending \(X \to \xi\) is a Lie algebra homomorphism. Reciprocally, given \(\xi\in \mathscr{C}^\infty(P; \mathfrak{g})\), we define \(X\in \Gamma(M; \mathrm{ad(P)})\) by setting \(X(m)=[p, \xi(p)]\) for any \(p \in \pi^{-1}(m)\). This is well-defined, and it is smooth because the following diagram is commutative
%     \begin{center}
%     \begin{tikzcd}[row sep=normal, column sep=normal, text height=1ex, text depth=1ex, ampersand replacement=\&]
%             P \arrow[r, "{\displaystyle \mathrm{id}_P \times \xi}"] \arrow[d, "{\displaystyle \pi}", '] \& P\times \mathfrak{g} \arrow[d, "{\displaystyle q}"]  \\
%              M  \arrow[r, "{\displaystyle X}"]\& P\times_G\mathfrak{g} 
%     \end{tikzcd} 
%     \end{center}
%     and we can use with \(\pi\) the characteristic property of surjective smooth surjections: \(X \circ \pi= q \circ (\mathrm{id}_P \times \xi), \) the latter being smooth. Finally, it is clear that these morphisms are inverses of each other. \\  
    
%     Let us see that \((\mathtt{G}, [\cdot, \cdot])\cong (\mathfrak{G}, [\cdot, \cdot])\). To do so, recall the definition of the Killing vector field associated to an element of a Lie algebra (equation \ref{def_Killing}). We define \(\sigma\colon \mathtt{G} \to \mathfrak{G}, \xi \to \sigma(\xi)\) as follows
%     \[
%     \sigma(\xi)_p:=X_{\xi(p)}=\left.\frac{\diff}{\diff t} \right|_{t=0} p \cdot \exp(t \cdot \xi(p)). 
%     \]
%     Then, {\color{red} \(\sigma\) is a Lie algebra isomorphism}. The vector field \(\sigma(\xi)\) is vertical because \(\pi\) is constant on \(G\)-orbits and
%     \[
%     \diff_p\pi (\sigma(\xi)_p)=\diff_p\pi \parent{\left.\frac{\diff}{\diff t} \right|_{t=0} p \cdot \exp(t \cdot \xi(p))} =\left.\frac{\diff}{\diff t}  \right|_{t=0} \pi \big(p \cdot \exp(t \cdot \xi(p)) \big) = 0.
%     \]
%     Moreover, \((\phi_g)_*\sigma(\xi)=\sigma(\xi)\). Certainly,
%     \begin{align*}
%         \diff_p\phi_g\parent{\sigma(\xi)_p} & = \diff_p \phi_g \parent{\left.\frac{\diff}{\diff t} \right|_{t=0} p \cdot \exp(t \cdot \xi(p))} = \left.\frac{\diff}{\diff t} \right|_{t=0} p \cdot \exp(t \cdot \xi(p)) \ g
%         =\left.\frac{\diff}{\diff t} \right|_{t=0} p \cdot \exp(t \cdot \xi(p\cdot g\ g^{-1})) \\
%         & = \left.\frac{\diff}{\diff t} \right|_{t=0} p \cdot \exp(t \cdot \Ad{g}\xi(p)) \ g = \left.\frac{\diff}{\diff t} \right|_{t=0} p \cdot g \ \exp(t \cdot \xi(p\cdot g)) \ g^{-1} \ g \\
%         & = \left.\frac{\diff}{\diff t} \right|_{t=0} (p\cdot g) \cdot \exp(t \cdot \xi(p\cdot g)) \\
%         & = \sigma(\xi)_{p \cdot g}.
%     \end{align*}
%     Thus, \(\sigma(\xi) \in \mathfrak{G}\). Moreover, \(\sigma\) is a Lie algebra homomorphism. {\color{red} A partir de aquí no lo tengo nada claro, revisar} To justify that, let us first observe that, for each \(p_0 \in P\), the restriction \(\sigma(\xi)|_{p_0 \cdot G}\) is a vector field on \(p_0 \cdot G\) (because \(\sigma(\xi)\) is vertical). Moreover,
%     \begin{align*}
%         \sigma(\xi)_{p_0 \cdot g} & = \left.\frac{\diff}{\diff t} \right|_{t=0} (p_0\cdot g) \cdot \exp(t \cdot \xi(p_0 \cdot g)) = \left.\frac{\diff}{\diff t} \right|_{t=0} (p\cdot g) \cdot \exp(t \cdot \Ad{g^{-1}}\xi(p_0)) \\
%         & =\left.\frac{\diff}{\diff t} \right|_{t=0} p_0\cdot \parent{g \ g^{-1} \cdot \exp(t \cdot \xi(p_0))\ g}= \left.\frac{\diff}{\diff t} \right|_{t=0} p_0 \cdot \exp(t \cdot \xi(p_0))\ g.
%     \end{align*}
%     Let \(\phi^{(p_0)}\colon G \to P, g \mapsto p_0 \cdot g\), and observe that, since \(\xi(p_0)\in \mathfrak{g}\), it is left-invariant, hence,
%     \[
%     \left.\frac{\diff}{\diff t} \right|_{t=0} \exp(t \cdot \xi(p_0)) \cdot g = \diff_e L_g (\xi(p_0)_e)=\xi(p_0)_g
%     \] 
%     where we used that \(\gamma=\exp(t\cdot \xi(p_0))\) is a smooth curve on \(G\) satisfying \(\gamma(0)=e\) and \(\gamma'(0)=\xi(p_0)_e\) (see \parencite[Cor. 3.25]{lee_introSM}). This implies that \(\sigma(\xi)_{p_0 \cdot g}=\diff_g\phi^{(p_0)}(\xi(p_0)_g)\), and therefore, \(\xi(p_0)\) and \(\sigma(\xi)|_{p_0 \cdot G}\) are \(\phi^{(p_0)}\)-related. Now, since \(p_0\cdot G\) is an embedded submanifold of \(P\), and \(\sigma(\xi)\) and \(\sigma(\eta)\) are tangent to that submanifold, \([\sigma(\xi), \sigma(\eta)]\) is tangent as well, and \([\sigma(\xi), \sigma(\eta)]_{p_0 \cdot g} = [\sigma(\xi)|_{p_0 \cdot G}, \sigma(\eta)|_{p_0 \cdot G}]_{p_0\cdot g}\). Therefore,
%     \begin{align*}
%     [\sigma(\xi), \sigma(\eta)]_{p_0 \cdot g} &= [\sigma(\xi)|_{p_0 \cdot G}, \sigma(\eta)|_{p_0 \cdot G}]_{p_0\cdot g} =[\phi^{(p_0)}_*(\xi(p_0)), \phi^{(p_0)}_*(\eta(p_0))]_{p_0 \cdot g}  \\
%     &= \phi^{(p_0)}_* \big([\xi(p_0), \eta(p_0)]\big)_{p_0 \cdot g}=\diff_{g}\phi^{(p_0)}\big([\xi, \eta](p_0)\big) = \sigma([\xi, \eta])_{p_0\cdot g}.
%     \end{align*}
%     We conclude by observing that \(P\) is the union of \(G\)-orbits. \\

%     To obtain an inverse, given \(X \in \mathfrak{G}\), each \(X_p\) is a vertical, and \(\diff_e\phi^{(p)}\colon \mathfrak{g} \to V_pP \subset T_pP\) is a vector space isomorphism. Hence, for every \(G\)-invariant vector field, we may define \(\xi(p):=\parent{\diff_e\phi^{(p)}}^{-1}(X_p)\). This defines a smooth map \(\xi\colon P \to \mathfrak{g}\). Similar arguments to the ones already given show that \(\xi\) is smooth and well-defined, and this map is again a lie algebra homomorphism that is the inverse of \(\sigma\).
% \end{proof}

